\documentclass[11pt]{article}

\usepackage[T1]{fontenc}
\usepackage{lmodern}
\usepackage[a4paper,margin=30mm]{geometry}
\setlength{\emergencystretch}{2em}
\usepackage{amsmath,amssymb,amsthm,mathtools}
\usepackage{microtype}
\usepackage{xcolor}
\usepackage[hypertexnames=false,colorlinks=true,linkcolor=blue!45!black,
  citecolor=blue!45!black,urlcolor=blue!55!black]{hyperref}

\newtheorem{theorem}{Theorem}[section]
\newtheorem{proposition}[theorem]{Proposition}
\newtheorem{lemma}[theorem]{Lemma}
\newtheorem{corollary}[theorem]{Corollary}
\theoremstyle{remark}
\newtheorem{remark}[theorem]{Remark}

\newcommand{\C}{\mathbb C}
\newcommand{\R}{\mathbb R}
\newcommand{\N}{\mathbb N}
\newcommand{\Rea}{\operatorname{Re}}
\newcommand{\Mell}{\mathcal M}
\newcommand{\GRH}{\mathrm{GRH}}
\newcommand{\RH}{\mathrm{RH}}
\newcommand{\eps}{\varepsilon}

\title{One-Sided Additive Convolution Laws for Dirichlet \(L\)-Functions:\\
Zero Abscissae, Parity Gram Identities, and Euler-Prime Sign Walls}
\author{Michael Haufschild}
\date{}

\begin{document}
\maketitle

\begin{abstract}
Let $\chi$ be a primitive Dirichlet character, let
$M_\chi(N)=\sum_{n\leq N}\mu(n)\chi(n)$, and let $\Theta_\chi$ be the supremum of
the real parts of the nontrivial zeros of $L(s,\chi)$.  For $k\geq1$, form the real
mixed additive convolution
\[
 C_{\chi,k}(N)=
 \sum_{\substack{a_1+\cdots+a_k+b_1+\cdots+b_k=N\\a_j,b_j\geq0}}
 \prod_{j=1}^k M_\chi(a_j)
 \prod_{j=1}^k \overline{M_\chi(b_j)},
 \qquad M_\chi(0)=0.
\]
We prove the exact one-sided order law
\[
 \limsup_{N\to\infty}
 \frac{\log\bigl(1+[C_{\chi,k}(N)]_+\bigr)}{\log(N+1)}
 =2k(1+\Theta_\chi)-1.
\]
Thus the positive excursions alone, with every negative coefficient discarded, recover the
rightmost zero abscissa at every even nonlinear order.  In particular, $\GRH$ for
$L(s,\chi)$ is equivalent to the value $3k-1$ at any one order $k$.  Independently of
$\GRH$, functional-equation symmetry yields
$C_{\chi,k}(N)=\Omega_+(N^{3k-1-\eps})$ for every $\eps>0$; hence every primitive
character has near-critical positive excursions at every mixed order.  At the raw coefficient
level, if $R_{\chi,k}$ is the Hermitian $2k$-fold additive convolution of
$\mu(n)\chi(n)$, we prove the stronger even-index law
$\limsup \log(1+[R_{\chi,k}(2N)]_+)/\log(N+1)\geq2k\Theta_\chi-1$ and show that
$R_{\chi,k}(2N)>0$ infinitely often.  Thus
$\sum_{a+b=2N}\mu(a)\mu(b)>0$ for infinitely many $N$, unconditionally.  More generally, the
even binary convolution of every completely multiplicative sequence satisfies an exact dyadic
Gram identity.  For the Liouville function this identity, together with the structural pole of
$\zeta(2s)/\zeta(s)$ at $s=1/2$, proves that same-sign pairs form a strict majority among the
ordered decompositions of infinitely many even integers.  For the integrated M\"obius
antidiagonal, we also prove an exact local-factor classification: removal of one Euler prime gives
a nonnegative discrete-tent kernel, prime $2$ is pointwise optimal, and removal of any second
prime necessarily creates a negative transfer lobe.  Thus coefficientwise local-factor positivity
has depth exactly one.  The proof combines an
exact Hermitian Abel identity with the Mellin transform
$\Gamma(s)/L(s,\chi)$ and the classical twisted Littlewood zero-abscissa theorem.  We also
state an abstract one-sided spectral law for reciprocal Dirichlet series and record a Lean~4
formalization of its Hermitian convolution and Abel core.  The theorem does not prove $\GRH$:
establishing the critical upper value remains precisely the unresolved source problem.
\end{abstract}

\medskip
\noindent\textbf{Keywords.} Dirichlet $L$-function; Mertens function; M\"obius function;
Liouville function; additive convolution; discrete tent; parity projection; Mellin transform; generalized
Riemann hypothesis.

\medskip
\noindent\textbf{2020 Mathematics Subject Classification.} 11M26, 11M06, 11N37, 40E05.

\section{Introduction}

The classical Littlewood criterion reads the zero-free abscissa of the Riemann zeta function
from the growth of the Mertens function.  Its fixed-character form says that the rightmost zero
of $L(s,\chi)$ determines the order of
\[
 M_\chi(x)=\sum_{n\leq x}\mu(n)\chi(n).
\]
This is a linear statement with an absolute value.  We show that the same zero coordinate is
encoded, without loss, by a family of nonlinear additive statistics in which all negative
excursions are ignored.

Write $[y]_+=\max(y,0)$.  The main theorem states that for every primitive character and every
$k\geq1$,
\[
 \Lambda_{\chi,k}:=\limsup_{N\to\infty}
 \frac{\log(1+[C_{\chi,k}(N)]_+)}{\log(N+1)}
 =2k(1+\Theta_\chi)-1.
\]
The affine dependence on $k$ is rigid.  Raising the convolution order neither weakens nor
amplifies the zero obstruction: every order records the same $\Theta_\chi$.  The use of the
positive part is essential.  The theorem is not an estimate for $|C_{\chi,k}(N)|$ and makes no
claim about its negative excursions.

For complex characters, the correct object is Hermitian rather than a plain even power.  If
\[
 A_\chi(z)=\sum_{n\geq1}\mu(n)\chi(n)z^n,
 \qquad
 B_\chi(z)=\sum_{N\geq0}M_\chi(N)z^N,
\]
then $A_\chi(z)=(1-z)B_\chi(z)$.  For real $0<z<1$, the mixed convolution is defined so that
\[
 |A_\chi(z)|^{2k}
 =(1-z)^{2k}\sum_{N\geq0}C_{\chi,k}(N)z^N.
\]
The left side is nonnegative, so the full series is bounded above by the series of positive
coefficient parts.  Any improvement in their exponent would force excessive decay of
$A_\chi(e^{-x})$ as $x\downarrow0$.  Mellin transformation would then continue
$1/L(s,\chi)$ holomorphically across an actual zero, which is impossible.

This argument also gives an unconditional positive-side theorem.  The functional equation and
coefficient conjugation make the zero set of a fixed $L(s,\chi)$ symmetric under
$\rho\mapsto1-\overline\rho$, so $\Theta_\chi\geq1/2$.  Consequently, for every $\eps>0$,
infinitely many $N$ satisfy
\[
 C_{\chi,k}(N)\geq N^{3k-1-\eps}.
\]
At $k=1$, this is a near-quadratic $\Omega_+$ theorem for the reflected twisted-Mertens
convolution.  It is substantially different from current work on raw binary sums of the form
$\sum_{n<N}\mu(n)\mu(N-n)$.  Binary raw Möbius analogues remain out of reach in the setting of
Banks and Shparlinski~\cite{BanksShparlinski}; Mangerel's theorem for the Liouville function
establishes strict nonmaximality rather than a critical positive-side order
\cite{Mangerel}.  The two summations hidden in $M_\chi$ are exactly what permit the present
Mellin--Abel argument.

The raw binary problem nevertheless has a rigorous one-sided consequence.  Averaging its
generating function at $r$ and $-r$ extracts the even output indices and remains an average of two
Hermitian squares.  A finite positive tail on this single arithmetic progression would bound
$A_\chi(r)$ as $r\uparrow1$, and Mellin transformation would make $1/L(s,\chi)$ holomorphic
throughout $\Rea s>0$.  Nontrivial zeros rule this out.  We obtain unconditional positive
recurrence on even totals at every order and a power-law $\Omega_+$ theorem for every $k\geq2$.
For completely multiplicative coefficients, the local factor at $2$ converts the parity average
into a sum of two adjacent dyadic squares.  This is a sign theorem and an exact identity, not a
binary Chowla estimate.

Additive convolutions of arithmetic functions have a large literature, but the nearby results
address different questions.  Everlove proves meromorphic continuation for Dirichlet series
attached to additive convolutions when one ordinary generating function is meromorphic near
$z=1$, and separately analyzes $\sum M(n)n^{-s}$ and character partial sums
\cite{Everlove}.  Misra, Saha and Sharma study divisor convolutions and functions with absolutely
convergent Ramanujan expansions~\cite{MisraSahaSharma}.  Liflandsky gives a recent explicit
formula and a direct decay criterion for the discrete Möbius Laplace transform
\cite{Liflandsky}.  Cantarini, Gambini and Zaccagnini derive weighted explicit formulas for raw
Liouville and Möbius convolutions under RH and simplicity assumptions
\cite{CantariniGambiniZaccagnini}; their results do not include the unconditional positive
recurrence on even totals proved below.  Mangerel proves the strict-nonmaximality statement
$|\sum_{n<N}\lambda(n)\lambda(N-n)|<N-1$ for every $N\geq11$, but no aggregate sign conclusion
\cite{Mangerel}.  Kaczorowski and Perelli's general $\Omega$-theorem concerns Dirichlet
coefficients of polynomial combinations of Selberg-class $L$-functions
\cite{KaczorowskiPerelli}; the coefficients below are ordinary-additive Cauchy convolutions of
reciprocal coefficients.  None of these results states the positive-part order laws proved here.

Section~\ref{sec:abstract} proves an abstract one-sided spectral law.  Section~\ref{sec:dirichlet}
checks its hypotheses for primitive Dirichlet $L$-functions and derives the $\GRH$ and
$\Omega_+$ corollaries.  Section~\ref{sec:raw} treats raw Möbius convolutions, even-index
recurrence, and the dyadic Liouville Gram law.  Section~\ref{sec:boundary} distinguishes the result from neighboring
convolution problems, proves the Euler-prime tent/sign-wall theorem, and isolates the remaining
open source.  Section~\ref{sec:lean} records the
formal certificate for the Hermitian Abel mechanism.

\section{An abstract one-sided spectral law}
\label{sec:abstract}

Let $(a_n)_{n\geq1}$ be a complex sequence, set $S(0)=0$ and
\[
 S(N)=\sum_{n\leq N}a_n,
 \qquad
 A(z)=\sum_{n\geq1}a_nz^n,
 \qquad
 B(z)=\sum_{N\geq0}S(N)z^N.
\]
We assume polynomial growth whenever needed, so all power series below converge for $|z|<1$.
For $k\geq1$, define
\begin{equation}
 C_k(N)=
 \sum_{\substack{a_1+\cdots+a_k+b_1+\cdots+b_k=N\\a_j,b_j\geq0}}
 \prod_{j=1}^k S(a_j)\prod_{j=1}^k\overline{S(b_j)}.
 \tag{2.1}
\end{equation}

\begin{lemma}[Hermitian Abel identity]
For every $N$, the coefficient $C_k(N)$ is real.  For real $0<z<1$,
\begin{equation}
 |A(z)|^{2k}=(1-z)^{2k}\sum_{N\geq0}C_k(N)z^N.
 \tag{2.2}
\end{equation}
Consequently,
\begin{equation}
 |A(z)|^{2k}
 \leq(1-z)^{2k}\sum_{N\geq0}[C_k(N)]_+z^N.
 \tag{2.3}
\end{equation}
\end{lemma}

\begin{proof}
Exchanging the two $k$-tuples in (2.1) conjugates every summand and preserves the indexing set;
hence $\overline{C_k(N)}=C_k(N)$.  The partial-sum identity gives
$B(z)=A(z)/(1-z)$.  The coefficient of $z^N$ in
$B(z)^k\overline{B(z)}^k$ is (2.1), and for real $z$ the product equals $|B(z)|^{2k}$.
This proves (2.2).  Since the left side is nonnegative and
$C_k(N)\leq[C_k(N)]_+$ termwise, (2.3) follows.
\end{proof}

We isolate the analytic inputs.  Let $D(s)=\sum_{n\geq1}a_nn^{-s}$ in an initial right
half-plane.  Assume that:
\begin{enumerate}
\item for every $\eta>0$,
\begin{equation}
 |S(N)|\ll_\eta (N+1)^{\theta+\eta};
 \tag{2.4}
\end{equation}
\item $D$ has a meromorphic continuation and has poles $\rho$ with $\Rea\rho>0$ whose real
parts have supremum $\theta$;
\item for $\Rea s$ initially large,
\begin{equation}
 \int_0^\infty A(e^{-x})x^{s-1}\,dx=\Gamma(s)D(s).
 \tag{2.5}
\end{equation}
\end{enumerate}

\begin{theorem}[One-sided even-convolution spectral law]
\label{thm:abstract}
Under the three hypotheses above, for every $k\geq1$,
\begin{equation}
 \limsup_{N\to\infty}
 \frac{\log(1+[C_k(N)]_+)}{\log(N+1)}
 =2k(1+\theta)-1.
 \tag{2.6}
\end{equation}
\end{theorem}

\begin{proof}
Fix $\eta>0$.  There are at most $(N+1)^{2k-1}$ weak compositions of $N$ into $2k$
coordinates.  By (2.4), every product in (2.1) is
$O_\eta((N+1)^{2k(\theta+\eta)})$.  Therefore
\[
 |C_k(N)|\ll_\eta
 (N+1)^{2k(1+\theta)-1+2k\eta},
\]
which proves the upper inequality in (2.6).

Suppose the lower inequality fails.  Choose
\[
 \Lambda_k<\lambda<2k(1+\theta)-1,
 \qquad
 \Lambda_k:=\limsup_{N\to\infty}
 \frac{\log(1+[C_k(N)]_+)}{\log(N+1)},
\]
with $\lambda>0$.  Absorbing finitely many initial values into a constant gives
\begin{equation}
 [C_k(N)]_+\ll_\lambda(N+1)^\lambda.
 \tag{2.7}
\end{equation}
Put $z=e^{-x}$ with $0<x\leq1$.  The elementary bound
\[
 \sum_{N\geq0}(N+1)^\lambda e^{-Nx}\ll_\lambda x^{-\lambda-1}
\]
and $1-e^{-x}\leq x$ turn (2.3) into
\begin{equation}
 |A(e^{-x})|\ll x^{-b},
 \qquad
 b=\frac{\lambda+1}{2k}-1<\theta.
 \tag{2.8}
\end{equation}
The exponential decay as $x\to\infty$ and (2.8) show that
\[
 H(s):=\frac1{\Gamma(s)}\int_0^\infty A(e^{-x})x^{s-1}\,dx
\]
is holomorphic in $\Rea s>b$.  By (2.5), $H(s)=D(s)$ in the initial half-plane, so uniqueness
of analytic continuation makes $H$ a holomorphic continuation of $D$ throughout
$\Rea s>b$.  By the definition of $\theta$, however, $D$ has a pole $\rho$ with
$\Rea\rho>b$.  This contradiction proves the lower inequality.
\end{proof}

\begin{remark}
The lower argument is one-sided.  It never bounds $|C_k(N)|$ and never uses
$[{-C_k(N)}]_+$.  The nonnegative radial quantity $|A(e^{-x})|^{2k}$ forces positive
coefficient mass at the exact spectral order.
\end{remark}

\section{Primitive Dirichlet characters}
\label{sec:dirichlet}

Let $\chi$ be a primitive Dirichlet character of conductor $q$.  For $q=1$, we use the trivial
character and $L(s,\chi)=\zeta(s)$.  For $q>1$, $\chi$ is nonprincipal.  Euler products give
\begin{equation}
 \frac1{L(s,\chi)}=\sum_{n\geq1}\frac{\mu(n)\chi(n)}{n^s},
 \qquad \Rea s>1.
 \tag{3.1}
\end{equation}
Let
\[
 \Theta_\chi=\sup\{\Rea\rho:L(\rho,\chi)=0,\ 0<\Rea\rho<1\}.
\]

We use the following classical fixed-character form of Littlewood's theorem.  It follows by the
usual reciprocal-$L$ Perron argument; see the treatments of reciprocal zeta and Dirichlet
$L$-functions in Titchmarsh~\cite{Titchmarsh} and Davenport~\cite{Davenport}.

\begin{proposition}[Twisted Littlewood zero-abscissa theorem]
\label{prop:littlewood}
For every fixed primitive $\chi$ and every $\eta>0$,
\begin{equation}
 M_\chi(N)=\sum_{n\leq N}\mu(n)\chi(n)
 \ll_{\chi,\eta}(N+1)^{\Theta_\chi+\eta}.
 \tag{3.2}
\end{equation}
Equivalently, the abscissa of convergence of the reciprocal Dirichlet series (3.1) is
$\Theta_\chi$.
\end{proposition}

\begin{remark}
No uniformity in $q$ is asserted.  For fixed $\delta>0$, the completed functional equation,
finite order, and the logarithmic-derivative expansion over zeros give the standard Littlewood
bound for $1/L(s,\chi)$ in $\Rea s\geq\Theta_\chi+\delta$; Perron inversion yields (3.2).
Conversely, any smaller partial-sum exponent would, by Abel summation, continue (3.1)
holomorphically across a zero to its right.
\end{remark}

For $x>0$, termwise integration in $\Rea s>1$ gives
\begin{equation}
 \int_0^\infty
 \left(\sum_{n\geq1}\mu(n)\chi(n)e^{-nx}\right)x^{s-1}\,dx
 =\frac{\Gamma(s)}{L(s,\chi)}.
 \tag{3.3}
\end{equation}
The poles of $1/L(s,\chi)$ in the critical strip are exactly the nontrivial zeros of
$L(s,\chi)$, counted with multiplicity.  The hypotheses of Theorem~\ref{thm:abstract} therefore
hold with $a_n=\mu(n)\chi(n)$ and $\theta=\Theta_\chi$.

\begin{theorem}[Dirichlet-$L$ mixed additive order law]
\label{thm:dirichlet}
For every primitive character $\chi$ and every integer $k\geq1$,
\begin{equation}
 \boxed{\quad
 \Lambda_{\chi,k}=2k(1+\Theta_\chi)-1.
 \quad}
 \tag{3.4}
\end{equation}
\end{theorem}

\begin{proof}
Apply Theorem~\ref{thm:abstract} using Proposition~\ref{prop:littlewood} and (3.3).
\end{proof}

For complex $\chi$, the functional equation first exchanges $\chi$ and $\overline\chi$.
Combining it with coefficient conjugation shows that the zero set of the same $L(s,\chi)$ is
stable under $\rho\mapsto1-\overline\rho$.  Hence $\Theta_\chi\geq1/2$, and equality holds if
and only if every nontrivial zero lies on the critical line.

\begin{corollary}[One-order $\GRH$ criterion]
For every fixed primitive $\chi$ and every fixed $k\geq1$,
\[
 \GRH\text{ for }L(s,\chi)
 \quad\Longleftrightarrow\quad
 \Lambda_{\chi,k}=3k-1.
\]
Thus Dirichlet $\GRH$ is equivalent to this identity for all primitive characters at any one
fixed order $k$.
\end{corollary}

\begin{corollary}[Unconditional positive excursions]
\label{cor:omega}
For every primitive $\chi$, every $k\geq1$, and every $\eps>0$,
\begin{equation}
 C_{\chi,k}(N)=\Omega_+\bigl(N^{3k-1-\eps}\bigr).
 \tag{3.5}
\end{equation}
More explicitly, infinitely many $N$ satisfy
$C_{\chi,k}(N)\geq N^{3k-1-\eps}$.
\end{corollary}

\begin{proof}
The zero symmetry gives $\Theta_\chi\geq1/2$, so (3.4) gives
$\Lambda_{\chi,k}\geq3k-1$.  By the definition of the limsup, for every $\eps>0$ there are
infinitely many $N$ for which
$1+[C_{\chi,k}(N)]_+>(N+1)^{3k-1-\eps/2}$.  For sufficiently large such $N$, this implies
$C_{\chi,k}(N)\geq N^{3k-1-\eps}$.
\end{proof}

\begin{corollary}[Zeta specialization]
Let $M(N)=\sum_{n\leq N}\mu(n)$ and
\[
 C_k(N)=\sum_{n_1+\cdots+n_{2k}=N}\prod_{j=1}^{2k}M(n_j).
\]
If $\Theta$ is the supremum of the real parts of the nontrivial zeta zeros, then
\[
 \limsup_{N\to\infty}\frac{\log(1+[C_k(N)]_+)}{\log(N+1)}
 =2k(1+\Theta)-1.
\]
In particular, $\RH$ is equivalent to the value $3k-1$, and for every $\eps>0$,
unconditionally,
\[
C_k(N)=\Omega_+(N^{3k-1-\eps}).
\]
\end{corollary}

\section{Raw convolutions and positive recurrence}
\label{sec:raw}

The same radial mechanism applies before taking partial sums.  Put
$a_\chi(n)=\mu(n)\chi(n)$ and, for $k\geq1$, define
\begin{equation}
 R_{\chi,k}(N)=
 \sum_{\substack{n_1+\cdots+n_k+m_1+\cdots+m_k=N\\n_j,m_j\geq1}}
 \prod_{j=1}^k a_\chi(n_j)
 \prod_{j=1}^k\overline{a_\chi(m_j)}.
 \tag{4.1}
\end{equation}
Exchanging the two tuples shows that $R_{\chi,k}(N)$ is real.  Absolute convergence gives the
exact radial identity
\begin{equation}
 \sum_{N\geq1}R_{\chi,k}(N)r^N=|A_\chi(r)|^{2k}\geq0,
 \qquad 0<r<1.
 \tag{4.2}
\end{equation}
Define
\begin{equation}
 \Psi_{\chi,k}=\limsup_{N\to\infty}
 \frac{\log(1+[R_{\chi,k}(N)]_+)}{\log(N+1)}.
 \tag{4.3}
\end{equation}

\begin{theorem}[Raw positive-excursion law]
\label{thm:raw}
For every primitive $\chi$ and every $k\geq1$,
\begin{equation}
 2k\Theta_\chi-1\leq\Psi_{\chi,k}\leq2k-1.
 \tag{4.4}
\end{equation}
Moreover, $R_{\chi,k}(N)>0$ for infinitely many $N$.  Consequently, for $k\geq2$ and every
$\eps>0$,
\begin{equation}
 R_{\chi,k}(N)=\Omega_+\bigl(N^{k-1-\eps}\bigr).
 \tag{4.5}
\end{equation}
\end{theorem}

\begin{proof}
There are at most $(N+1)^{2k-1}$ compositions in (4.1), and every summand has modulus at most
one.  This proves the upper inequality in (4.4).

For the lower inequality, the assertion is automatic if $2k\Theta_\chi-1\leq0$.  Otherwise,
suppose that $\Psi_{\chi,k}<\lambda<2k\Theta_\chi-1$.  Then
$[R_{\chi,k}(N)]_+\ll_\lambda(N+1)^\lambda$.  With $r=e^{-x}$, (4.2) gives
\[
 |A_\chi(e^{-x})|^{2k}
 \leq\sum_{N\geq1}[R_{\chi,k}(N)]_+e^{-Nx}
 \ll_\lambda x^{-\lambda-1}.
\]
Thus $A_\chi(e^{-x})=O(x^{-b})$ for
$b=(\lambda+1)/(2k)<\Theta_\chi$.  As in the proof of
Theorem~\ref{thm:abstract}, Mellin transformation makes $1/L(s,\chi)$ holomorphic in
$\Rea s>b$, contradicting a nontrivial zero to the right of $b$.  This proves (4.4).

It remains to handle positive recurrence, including the endpoint $k=1$ when the lower exponent
can be zero.  If only finitely many $R_{\chi,k}(N)$ were positive, (4.2) would give a uniform
bound for $|A_\chi(r)|^{2k}$ on $0<r<1$.  Hence $A_\chi(e^{-x})$ would be bounded at zero and
exponentially decaying at infinity.  Its Mellin transform would continue $1/L(s,\chi)$
holomorphically throughout $\Rea s>0$, contradicting the existence of nontrivial zeros.
Therefore positive coefficients occur infinitely often.  Finally,
$\Theta_\chi\geq1/2$ and (4.4) give $\Psi_{\chi,k}\geq k-1>0$ for $k\geq2$, which is (4.5).
\end{proof}

The unrestricted theorem admits a strict arithmetic-progression refinement.  Define
\[
 \Psi^{\mathrm{ev}}_{\chi,k}=\limsup_{q\to\infty}
 \frac{\log(1+[R_{\chi,k}(2q)]_+)}{\log(q+1)}.
\]

\begin{theorem}[Even-index Hermitian pole law]
\label{thm:parity}
For every primitive $\chi$ and every $k\geq1$,
\begin{equation}
 \sum_{q\geq0}R_{\chi,k}(2q)u^q
 =\frac{|A_\chi(\sqrt u)|^{2k}+|A_\chi(-\sqrt u)|^{2k}}2\geq0,
 \qquad 0<u<1.                                                \tag{4.6}
\end{equation}
Moreover,
\begin{equation}
 2k\Theta_\chi-1\leq\Psi^{\mathrm{ev}}_{\chi,k}\leq2k-1,    \tag{4.7}
\end{equation}
and $R_{\chi,k}(2q)>0$ for infinitely many $q$.  For $k\geq2$ and every $\eps>0$,
$R_{\chi,k}(2q)\geq q^{k-1-\eps}$ for infinitely many $q$.
\end{theorem}

\begin{proof}
Let
\[
 F_{\chi,k}(z)=A_\chi(z)^k
 \left(\sum_{n\geq1}\overline{a_\chi(n)}z^n\right)^k.
\]
Its coefficient of $z^N$ is $R_{\chi,k}(N)$, while
$F_{\chi,k}(r)=|A_\chi(r)|^{2k}$ for real $r$.  The order-two character projector gives
\[
 \frac{F_{\chi,k}(r)+F_{\chi,k}(-r)}2
 =\sum_{q\geq0}R_{\chi,k}(2q)r^{2q},
\]
which is (4.6).  In particular,
\[
 |A_\chi(\sqrt u)|^{2k}
 \leq2\sum_{q\geq0}[R_{\chi,k}(2q)]_+u^q.                   \tag{4.8}
\]
If the lower inequality in (4.7) failed, choose
$\Psi^{\mathrm{ev}}_{\chi,k}<\lambda<2k\Theta_\chi-1$ and put
$u=e^{-2x}$ in (4.8).  The Abel moment bound would give
$A_\chi(e^{-x})=O(x^{-(\lambda+1)/(2k)})$, continuing $1/L(s,\chi)$
holomorphically across a zero to the right.  The upper inequality follows from the number of
compositions as before.

If only finitely many even-index coefficients were positive, the right side of (4.8) would be
bounded as $u\uparrow1$.  The same Mellin integral would then continue $1/L(s,\chi)$ throughout
$\Rea s>0$, again contradicting a nontrivial zero.  Finally $\Theta_\chi\geq1/2$ gives the stated
higher-order positive excursions.
\end{proof}

\begin{corollary}[Binary additive M\"obius positivity on even totals]
For infinitely many positive integers $N$,
\begin{equation}
 \sum_{a+b=2N}\mu(a)\mu(b)\geq1.                              \tag{4.9}
\end{equation}
\end{corollary}

\begin{proof}
Take the conductor-one character and $k=1$ in Theorem~\ref{thm:parity}.  The raw coefficients are
integers, so every positive coefficient is at least one.
\end{proof}

There is a further exact refinement for completely multiplicative sequences.  Let $f$ be
complex-valued and completely multiplicative, put $c=f(2)$ and
\[
 F_f(r)=\sum_{n\geq1}f(n)r^n,
 \qquad H_f(N)=\sum_{a+b=N}f(a)\overline{f(b)}.
\]
Since
\[
 F_f(r)+F_f(-r)=2\sum_{n\ \mathrm{even}}f(n)r^n=2cF_f(r^2),
\]
the binary parity projector and the parallelogram identity give the following universal law.

\begin{proposition}[Dyadic Gram law]
\label{prop:dyadic}
For every real $0<r<1$ for which the power series converges absolutely,
\begin{equation}
 \boxed{\quad
 \sum_{q\geq0}H_f(2q)r^{2q}
 =|F_f(r)-cF_f(r^2)|^2+|cF_f(r^2)|^2.
 \quad}                                                       \tag{4.10}
\end{equation}
\end{proposition}

More generally, full multiplicativity is not needed.  If a sequence satisfies $a(2n)=ca(n)$ and
$O(r)=\sum_{n\ \mathrm{odd}}a(n)r^n$, then even--odd splitting gives
\begin{equation}
 A(r)=O(r)+cA(r^2),
 \qquad
 A(r)=\sum_{j\geq0}c^jO(r^{2^j}).                             \tag{4.10a}
\end{equation}
The second identity follows by iteration; its remainder tends to zero doubly exponentially.
Conversely, within the polynomial-growth class, arbitrary odd-index data $O$ in (4.10a) give the
unique sequence satisfying the local recursion.  Proposition~\ref{prop:dyadic} is therefore the
orthogonal decomposition
\begin{equation}
 \sum_{q\geq0}H_f(2q)r^{2q}=|O(r)|^2+|cF_f(r^2)|^2.           \tag{4.10b}
\end{equation}
At the Dirichlet-series level,
\[
 D_{\mathrm{odd}}(s)=(1-c2^{-s})D(s).
\]
This classifies the exact freedom left by the dyadic law: every odd-index coefficient remains
source data.

For $f=\lambda$, $c=-1$ and
$\sum_n\lambda(n)n^{-s}=\zeta(2s)/\zeta(s)$.  The quotient has a structural pole at
$s=1/2$, since $\zeta(1/2)$ is finite and nonzero.  If the even coefficients in (4.10) were
positive only finitely often, its right side would bound $F_\lambda(r)$ at $r=1$ and the Mellin
transform would continue $\zeta(2s)/\zeta(s)$ holomorphically through $s=1/2$.  Hence:

\begin{corollary}[Liouville same-sign majority]
\label{cor:liouville}
For infinitely many positive integers $N$,
\begin{equation}
 \sum_{a+b=2N}\lambda(a)\lambda(b)\geq1.                     \tag{4.11}
\end{equation}
Equivalently, among the $2N-1$ ordered positive decompositions of these even integers, same-sign
Liouville pairs form a strict majority over opposite-sign pairs.
\end{corollary}

For Liouville, the multiplier $1+2^{-s}$ relating the odd seed to the full Dirichlet series has
all its zeros on $\Rea s=0$.  It is therefore an analytic unit on $\Rea s>0$: the dyadic Gram law
exposes two nonnegative channels but cannot remove any right-half pole.  An upper theorem for the
odd seed remains a genuine source problem, not a consequence of parity projection.

For $m\geq2$, the known generalized Liouville function $\lambda_m$ of de
Camargo~\cite{Camargo} has Dirichlet series $\zeta(ms)/\zeta(s)$ and local coefficients
$1,-1,0,\ldots,0,1,-1,\ldots$ of period $m$ on every prime-power tower.  The proof of
Theorem~\ref{thm:parity} applies verbatim: its Hermitian raw convolutions are positive at infinitely many even totals, and
their even-index positive-part exponent is at least $2k\Theta-1$.  The family interpolates
$\lambda_2=\lambda$ and the pointwise limit $\lambda_m\to\mu$; the family itself is not new.

Theorems~\ref{thm:raw} and~\ref{thm:parity} use no RH, zero simplicity, linear-independence, or zero-density
hypothesis.  It is weaker than the binary cancellation estimates sought in current work:
positive recurrence neither bounds $|R_{\chi,1}(N)|$ nor proves $R_{\chi,1}(N)=o(N)$.

\section{Boundary with raw binary correlations and the remaining source}
\label{sec:boundary}

Theorem~\ref{thm:dirichlet} does not estimate the raw binary sum
\[
 R_\chi(N)=\sum_{n=1}^{N-1}\mu(n)\chi(n)\mu(N-n)\chi(N-n).
\]
For $k=1$, the object in this paper is the twice-discrete-integrated correlation
\[
 C_{\chi,1}(N)=\sum_{a+b=N}M_\chi(a)\overline{M_\chi(b)}.
\]
Indeed, if $\Delta f(N)=f(N)-f(N-1)$, then the generating-function identity shows that the
appropriate second difference of $C_{\chi,1}$ recovers the corresponding raw additive
coefficient.  Two integrations are analytically decisive: they introduce the factor
$(1-z)^{-2}$ that shifts the coefficient order while retaining the Mellin singularity of
$1/L(s,\chi)$.

The upper value in the $\GRH$ corollary is not automatic.  At $k=1$, proving for every
$\eps>0$ that
\begin{equation}
 C_{\chi,1}(N)\ll_{\chi,\eps}N^{2+\eps}
 \tag{5.1}
\end{equation}
already implies $\GRH$ for $L(s,\chi)$.  Theorem~\ref{thm:dirichlet} sharpens this statement:
if an off-line zero exists, then positive values alone must violate (5.1) at the exact exponent
$2(1+\Theta_\chi)-1$.  Higher order cannot evade that obstruction, because all $k$ encode the
same $\Theta_\chi$.

The parity refinement does not evade it either.  Equations (4.10a)--(4.10b) classify every
sequence with the local rule $a(2n)=ca(n)$ by a free odd-index seed.  In the Liouville case the
factor relating that seed to the full Dirichlet series is $1+2^{-s}$, an analytic unit throughout
$\Rea s>0$.  Thus the dyadic Gram decomposition neither cancels nor moves a right-half pole.
It converts a hypothetical upper bound for even convolution coefficients into transform control,
but supplies no such coefficient bound itself.

There is nevertheless an exact local-factor law for the integrated M\"obius antidiagonal.  Put
$M(a)=\sum_{n\leq a}\mu(n)$, $U(N)=\sum_{a+b=N}M(a)M(b)$, and
$w_N(t)=(N-t+1)_+$.

\begin{proposition}[Euler-prime tent and two-prime sign wall]
\label{prop:prime-tent}
Let $p$ be prime.  Then
\begin{equation}
 M(a)=\sum_{\substack{p\nmid r\\r\leq a<pr}}\mu(r)                 \tag{5.2}
\end{equation}
and
\begin{equation}
 U(N)=\sum_{\substack{p\nmid r\\p\nmid s}}\mu(r)\mu(s)
 K_{p,N}(r,s),                                                    \tag{5.3}
\end{equation}
where
\begin{align}
 K_{p,N}(r,s)
 &=w_N(r+s)-w_N(pr+s)-w_N(r+ps)+w_N(p(r+s))                       \tag{5.4}\\
 &=A-(A-R)_+-(A-S)_++(A-R-S)_+,                                  \tag{5.5}
\end{align}
with $A=(N+1-r-s)_+$, $R=(p-1)r$, and $S=(p-1)s$.  In particular,
$K_{p,N}(r,s)$ is a nonnegative integer, and for $r,s\geq1$,
\begin{equation}
 K_{p,N}(r,s)>0
 \quad\Longleftrightarrow\quad
 r+s\leq N\ \text{ and }\ N+2\leq p(r+s).                       \tag{5.6}
\end{equation}
Moreover $K_{2,N}(r,s)\leq K_{p,N}(r,s)$ for every prime $p$.

In contrast, if $2\leq p<q$ and
\[
 Q_{p,q}(z)=\frac{z-z^p-z^q+z^{pq}}{1-z},
 \qquad k=\max(q+1,2p-1),
\]
then
\begin{equation}
 [z^k]Q_{p,q}(z)^2<0.                                             \tag{5.7}
\end{equation}
For $(p,q)=(2,3)$, the coefficient in degree $4$ is $-2$.
\end{proposition}

\begin{proof}
Remove the Euler factor at $p$ coefficientwise:
\[
 \mu(n)=\mu^{(p)}(n)-\mathbf 1_{p\mid n}\mu^{(p)}(n/p),
 \qquad
 \mu^{(p)}(r)=\mathbf 1_{p\nmid r}\mu(r).
\]
Summation gives (5.2), or equivalently the coefficientwise formal identity
\[
 \sum_{a\geq0}M(a)z^a
 =\sum_{p\nmid r}\mu(r)\frac{z^r-z^{pr}}{1-z}
 =\sum_{p\nmid r}\mu(r)\sum_{a=r}^{pr-1}z^a.
\]
Squaring and extracting degree $N$ proves (5.3)--(5.4).  After translating the two intervals by
$(r,s)$, the kernel counts
\[
 \#\{(i,j):0\leq i<R,\ 0\leq j<S,\ i+j=A-1\}.
\]
Inclusion--exclusion gives (5.5); nonnegativity, (5.6), and the bound by each of $R,S$ follow at
once.  Increasing $p$ enlarges both interval polynomials coefficientwise, proving the asserted
minimality of $p=2$.

For the sign wall, polynomial division gives
\[
 Q_{p,q}(z)=(z+\cdots+z^{p-1})-(z^q+\cdots+z^{pq-1}).
\]
The choice of $k$ satisfies $2p-2<k<2q$.  Thus its squared coefficient has no
positive--positive or negative--negative contribution.  A cross contribution exists because
$1\leq k-q\leq p-1$ after choosing the appropriate endpoint case.  Every surviving contribution
is therefore negative and occurs with its transposed partner.  For $(2,3,4)$ the two terms are
$(1,3)$ and $(3,1)$, giving $-2$.
\end{proof}

Proposition~\ref{prop:prime-tent} is a sharp local classification, not an upper estimate for
$U(N)$.  It says that coefficientwise local-factor positivity has depth exactly one: prime $2$
gives the smallest positive tent, while every attempted second removal necessarily creates a
negative lobe.  The remaining source must exploit signed multi-prime interactions.  Replacing the
odd M\"obius seed by an arbitrary bounded seed cannot work: the constant odd seed makes its window
sum grow linearly and its reflected antidiagonal cubically.

This identifies the actual frontier.  Neither ordinary $L^2$ control nor a phase average is
enough: such positive norms discard the signed reflected covariance that appears pointwise in
(5.1).  Current estimates for Möbius correlations do not supply the needed growing-reflection
upper law.  The contribution of the present theorem is therefore an exact spectral classification,
an unconditional raw sign-recurrence theorem, and positive-excursion laws, not a proof of
$\GRH$.

\section{Lean formalization}
\label{sec:lean}

The algebraic and Abel core has been formalized in Lean~4 against mathlib.  The module
\begin{center}
\small
\texttt{GaloisForLFunctions.RHTower.}\\
\texttt{DirichletLMixedAdditiveZeroAbscissaLaw}
\end{center}
defines the mixed additive convolution for an arbitrary complex sequence and proves:
\begin{enumerate}
\item every mixed coefficient is fixed by conjugation and has zero imaginary part;
\item absolute summability of the original Abel series implies absolute summability of every
mixed series;
\item for conjugation-fixed $q$,
\[
 \left(\sum_Na_Nq^N\right)^k
 \overline{\left(\sum_Na_Nq^N\right)^k}
 =\sum_N C_k(N)q^N;
\]
\item coefficientwise upper bounds alone imply the corresponding norm-square Abel bound, with
all negative coefficients discarded.
\end{enumerate}
The companion module
\begin{center}
\small
\texttt{GaloisForLFunctions.RHTower.}\\
\texttt{DirichletLMixedAdditiveEvenIndexLaw}
\end{center}
proves the exact order-two projector (4.6), nonnegativity of its real Abel sum, and the complex
parallelogram identity behind Proposition~\ref{prop:dyadic}.  It reuses the preceding Hermitian
Cauchy-power machinery; no duplicate convolution hierarchy is introduced.
The module
\begin{center}
\small
\texttt{GaloisForLFunctions.RHTower.}\\
\texttt{MobiusMertensPrimeRemovalTentLaw}
\end{center}
proves the coefficientwise Euler-prime difference, the exact prime-free interval formula for each
Mertens partial sum, its connector to the existing signed-antidiagonal RH criterion, discrete-tent
positivity and monotonicity, and the two-prime sign wall of
Proposition~\ref{prop:prime-tent}.
The source compiles with
\begin{center}
\small
\texttt{lake env lean}\\
\texttt{GaloisForLFunctions/RHTower/}\\
\texttt{DirichletLMixedAdditiveEvenIndexLaw.lean}.
\end{center}
An axiom audit reports only the standard logical principles
\texttt{propext}, \texttt{Classical.choice}, and \texttt{Quot.sound}.  The definitions of
Dirichlet characters, $\Theta_\chi$, Proposition~\ref{prop:littlewood}, and the exact limsup law
remain paper-level; no formal $\GRH$ claim is made.

\section{Concluding perspective}

The rightmost zero of a primitive Dirichlet $L$-function has an exact nonlinear additive
signature.  Every mixed even order gives the same coordinate, and the positive coefficients
alone carry all of it.  This produces both an all-character $\GRH$ criterion and an unconditional
$\Omega_+$ theorem at the critical order.  Before summation, the parity-projected radial mechanism
proves that raw Hermitian convolutions have recurrent positive coefficients on even totals.  In
the zeta case this gives unconditional even-index positivity of the binary additive M\"obius
self-convolution.  For completely multiplicative sequences, the same projector is an exact
two-square dyadic Gram law; the Liouville specialization forces a same-sign majority along an
infinite sequence of even totals.

For the integrated M\"obius statistic, Euler-prime removal has an exact and rigid sign geometry:
one prime produces a nonnegative interval-overlap tent, prime $2$ is optimal, and a second prime
forces a negative lobe.  This closes coefficientwise local positivity as a source mechanism and
locates the unresolved arithmetic content in genuinely global signed prime interactions.

The remaining open problem is now sharply separated from the transform machinery: prove the
critical upper exponent for one order without assuming zero information.  For zeta, that means
controlling the positive side of $\sum_{a+b=N}M(a)M(b)$ at order $N^{2+o(1)}$.  Such a theorem
would prove $\RH$; no argument in this paper is claimed to provide it.

\begin{thebibliography}{99}
\raggedright

\bibitem{BanksShparlinski}
W. D. Banks and I. E. Shparlinski,
\emph{Multiple sums with the M\"obius function},
Q. J. Math. 77 (2026), 109--127.

\bibitem{CantariniGambiniZaccagnini}
M. Cantarini, A. Gambini and A. Zaccagnini,
\emph{On the discrete convolution of the Liouville and M\"obius functions},
arXiv:2603.10241 (2026).

\bibitem{Camargo}
A. P. de Camargo,
\emph{Two generalizations of Liouville $\lambda$ function},
Notes Number Theory Discrete Math. 29 (2023), no.~1, 30--39.

\bibitem{Davenport}
H. Davenport,
\emph{Multiplicative Number Theory}, 3rd ed., revised by H. L. Montgomery,
Graduate Texts in Mathematics 74, Springer, 2000.

\bibitem{Everlove}
C. Everlove,
\emph{Meromorphic Dirichlet Series}, Ph.D. thesis, University of Michigan, 2018,
\url{https://deepblue.lib.umich.edu/items/7997c479-ffa7-4330-bdae-e408bd0a2d9a}.

\bibitem{HardyLittlewood}
G. H. Hardy and J. E. Littlewood,
\emph{Tauberian theorems concerning power series and Dirichlet's series whose coefficients are
positive}, Proc. London Math. Soc. (2) 13 (1914), 174--191.

\bibitem{KaczorowskiPerelli}
J. Kaczorowski and A. Perelli,
\emph{General $\Omega$-theorems for coefficients of $L$-functions},
Proc. Amer. Math. Soc. 143 (2015), 5139--5145.

\bibitem{Liflandsky}
S. Liflandsky,
\emph{Explicit formula for the discrete Laplace transform of the M\"obius function, related
special functions, and a criterion for the Riemann hypothesis},
arXiv:2607.09797 (2026).

\bibitem{Mangerel}
A. P. Mangerel,
\emph{On a Goldbach-type problem for the Liouville function},
Int. Math. Res. Not. IMRN 2024, no.~16, 11865--11877.

\bibitem{MisraSahaSharma}
B. Misra, B. Saha and A. Sharma,
\emph{On additive convolution sum of arithmetic functions and related questions},
arXiv:2502.08305, version 2 (2026).

\bibitem{Titchmarsh}
E. C. Titchmarsh,
\emph{The Theory of the Riemann Zeta-Function}, 2nd ed., revised by D. R. Heath-Brown,
Oxford University Press, 1986.

\end{thebibliography}

\end{document}
